\documentclass[11pt, letterpaper]{article}

\usepackage[margin=1in]{geometry}
\usepackage{hyperref}
\usepackage{listings}
\usepackage{xcolor}
\usepackage{booktabs}
\usepackage{enumitem}
\usepackage{graphicx}
\usepackage{fancyhdr}
\usepackage{titlesec}

% ── Colours and style ──────────────────────────────────────────────────
\definecolor{codebg}{RGB}{245,245,245}
\definecolor{codeframe}{RGB}{200,200,200}
\definecolor{keyword}{RGB}{0,100,170}
\definecolor{comment}{RGB}{100,100,100}
\definecolor{warnbg}{RGB}{255,245,230}

\lstset{
  basicstyle=\ttfamily\small,
  backgroundcolor=\color{codebg},
  frame=single,
  rulecolor=\color{codeframe},
  breaklines=true,
  breakatwhitespace=false,
  showstringspaces=false,
  tabsize=2,
  keywordstyle=\color{keyword},
  commentstyle=\color{comment}\itshape,
  columns=fullflexible,
  keepspaces=true,
  aboveskip=8pt,
  belowskip=8pt,
}

\hypersetup{
  colorlinks=true,
  linkcolor=blue!60!black,
  urlcolor=blue!60!black,
  citecolor=blue!60!black,
}

\pagestyle{fancy}
\fancyhf{}
\rhead{\small openpilot-hcc Setup Guide}
\lhead{\small UVA Link Lab}
\cfoot{\thepage}

\titleformat{\section}{\large\bfseries}{}{0em}{}[\titlerule]
\titleformat{\subsection}{\normalsize\bfseries}{}{0em}{}

% ── Convenience commands ───────────────────────────────────────────────
\newcommand{\warn}[1]{%
  \medskip\noindent\colorbox{warnbg}{\parbox{\dimexpr\linewidth-2\fboxsep}{%
    \textbf{Note:} #1}}\medskip}

\newcommand{\repo}{\texttt{openpilot-hcc}}
\newcommand{\prefix}[1]{\texttt{OPENPILOT\_PREFIX=#1}}

\begin{document}

% ── Title ──────────────────────────────────────────────────────────────
\begin{center}
  {\LARGE\bfseries openpilot-hcc Setup Guide} \\[6pt]
  {\large Three Operational Modes: Radar-Only, V2V Simulation, and Hardware-in-the-Loop} \\[12pt]
  UVA Link Lab \quad|\quad \today
\end{center}

\tableofcontents
\newpage

% ======================================================================
\section{Overview}
% ======================================================================

The \repo{} fork of comma.ai's openpilot implements Human-in-the-Loop
Cooperative Cruise Control (hCCC) with three operational modes, each
building on the one before it:

\begin{description}[leftmargin=2em, labelwidth=2em]
  \item[Mode 1] \textbf{Single-device, radar-only.}
    One openpilot instance on the PC follows a simulated lead vehicle
    using radar perception. No V2V, no physical hardware.
  \item[Mode 2] \textbf{Dual-device V2V simulation.}
    Two openpilot instances on the PC (ego + lead) communicate through a
    local UDP V2V relay.  Both use the MetaDrive simulator running a
    single bridge process.
  \item[Mode 3] \textbf{Hardware-in-the-loop (HIL).}
    Two Comma~3X devices run the full on-device openpilot stack.  The PC
    runs MetaDrive as the simulated world, shipping camera frames and CAN
    data to each device over USB-RNDIS.  V2V flows over WiFi between the
    devices directly.
\end{description}

\subsection{Architecture Summary}

\begin{center}
\begin{tabular}{@{}llll@{}}
\toprule
  & \textbf{Mode 1} & \textbf{Mode 2} & \textbf{Mode 3} \\
\midrule
openpilot runs on & PC & PC ($\times$2) & Comma 3X ($\times$2) \\
MetaDrive runs on & PC & PC & PC \\
Lead perception   & Radar only & V2V (UDP relay) & V2V (WiFi) \\
Physical devices  & None & None & 2$\times$ Comma 3X + USB-C \\
Terminals needed  & 2 & 4 & 3 \\
\bottomrule
\end{tabular}
\end{center}


% ======================================================================
\section{Prerequisites (All Modes)}
\label{sec:prereqs}
% ======================================================================

\subsection{Host Machine}

\begin{itemize}
  \item \textbf{OS:} Ubuntu 22.04+ (recommended) or macOS 13+.
  \item \textbf{Python:} 3.11 or 3.12 (required by \texttt{pyproject.toml}).
  \item \textbf{GPU:} A discrete GPU is recommended for MetaDrive rendering
        but not strictly required.
\end{itemize}

\subsection{Clone and Build}

\begin{lstlisting}[language=bash]
# Clone the repository
git clone https://github.com/ethanmathias/openpilot-hcc.git
cd openpilot-hcc
git checkout hcc-ego

# Initialize submodules
git submodule update --init --recursive

# Set up the virtual environment and build
tools/op.sh setup
source .venv/bin/activate
scons -u -j$(nproc)
\end{lstlisting}

\warn{If \texttt{git lfs} objects are missing (files show as 130-byte
pointers), run \texttt{git lfs pull} before building.}

\subsection{Python Dependencies for Simulation}

The MetaDrive simulator and supporting packages are declared in
\texttt{pyproject.toml} under the \texttt{[tools]} extra:

\begin{lstlisting}[language=bash]
pip install -e ".[tools,dev]"
\end{lstlisting}

Key packages installed by this step:
\begin{itemize}
  \item \texttt{metadrive-simulator} --- the driving simulator
  \item \texttt{pygame} --- used for the HIL composite window and
        Logitech wheel input
  \item \texttt{av} (PyAV) --- H.264 encoding/decoding for HIL camera
        streaming
\end{itemize}


% ======================================================================
\section{Mode 1: Single-Device Radar-Only}
\label{sec:mode1}
% ======================================================================

This is the simplest mode.  One openpilot instance runs on the PC under
the \texttt{hccego} prefix.  MetaDrive spawns a lead vehicle driven by a
scenario profile (or IDM policy), and openpilot's hCCC controller tracks
it using simulated radar returns.

\subsection{Terminal Layout}

You need \textbf{two terminals}, both in the repo root.

\subsubsection{Terminal 1 --- openpilot}

\begin{lstlisting}[language=bash]
cd ~/openpilot-hcc
source .venv/bin/activate
bash tools/sim/launch_openpilot.sh
\end{lstlisting}

This script exports the required environment variables automatically:

\begin{center}
\begin{tabular}{@{}ll@{}}
\toprule
\textbf{Variable} & \textbf{Value} \\
\midrule
\texttt{SIMULATION}         & \texttt{1} \\
\texttt{NOBOARD}            & \texttt{1} \\
\texttt{SKIP\_FW\_QUERY}    & \texttt{1} \\
\texttt{OPENPILOT\_PREFIX}  & \texttt{hccego} \\
\texttt{FINGERPRINT}        & \texttt{HONDA\_CIVIC\_2022} \\
\texttt{BLOCK}              & camerad, loggerd, encoderd, micd, logmessaged \\
\bottomrule
\end{tabular}
\end{center}

It also writes the necessary params (\texttt{AlphaLongitudinalEnabled},
\texttt{EnableHCCC}) and then starts the openpilot manager.

\subsubsection{Terminal 2 --- MetaDrive bridge}

Wait until Terminal~1 shows openpilot processes starting (roughly 5--10
seconds), then:

\begin{lstlisting}[language=bash]
cd ~/openpilot-hcc
source .venv/bin/activate
OPENPILOT_PREFIX=hccego python tools/sim/run_bridge.py --mode hc3
\end{lstlisting}

\subsection{Expected Behavior}

\begin{itemize}
  \item The MetaDrive window opens showing the ego vehicle behind a lead.
  \item After a few seconds the system auto-engages (cruise buttons are
        cycled automatically).
  \item The ego vehicle follows the lead, decelerating and accelerating
        to match its speed.
  \item Terminal~1 shows openpilot logs; Terminal~2 prints periodic
        status lines with \texttt{vEgo}, \texttt{dRel}, and engagement
        state.
\end{itemize}

\subsection{Scenario Replay (Optional)}

To replay a specific lead speed profile from the scenarios CSV:

\begin{lstlisting}[language=bash]
OPENPILOT_PREFIX=hccego python tools/sim/run_bridge.py \
    --mode hc3 --scn 48
\end{lstlisting}

This uses the profile at column index~48 in
\texttt{tools/sim/lib/Scenarios.csv}.  When the profile completes, a
telemetry CSV and PNG speed graph are saved under
\texttt{tools/sim/data/} and \texttt{tools/sim/graphs/}.

\subsection{Keyboard Controls}

\begin{center}
\begin{tabular}{@{}ll@{}}
\toprule
\textbf{Key} & \textbf{Action} \\
\midrule
\texttt{W / S}           & Throttle / brake \\
\texttt{A / D}           & Steer left / right \\
\texttt{1}               & Cruise: engage (DECEL\_SET) \\
\texttt{2}               & Cruise: resume/accel \\
\texttt{3}               & Cruise: cancel \\
\texttt{4}               & Cruise: main toggle \\
\texttt{R}               & Reset scenario \\
\texttt{Q} or \texttt{Esc} & Quit \\
\bottomrule
\end{tabular}
\end{center}


% ======================================================================
\section{Mode 2: Dual-Device V2V Simulation}
\label{sec:mode2}
% ======================================================================

Two openpilot instances run on the same PC under separate
\texttt{OPENPILOT\_PREFIX} namespaces (\texttt{hccego} and
\texttt{hcclead}).  A local UDP V2V relay forwards lead state to the ego
so hCCC can use V2V data instead of (or alongside) radar.

\subsection{Terminal Layout}

You need \textbf{four terminals}, all in the repo root.

\subsubsection{Terminal 1 --- V2V Relay Server}

\begin{lstlisting}[language=bash]
cd ~/openpilot-hcc
source .venv/bin/activate
python tools/hcc_v2v/relay_server.py --host 127.0.0.1 --port 19090
\end{lstlisting}

Leave this running for the entire session.

\subsubsection{Terminal 2 --- Lead openpilot}

\begin{lstlisting}[language=bash]
cd ~/openpilot-hcc
source .venv/bin/activate

export SIMULATION=1 NOBOARD=1 SKIP_FW_QUERY=1
export OPENPILOT_PREFIX=hcclead
export FINGERPRINT=HONDA_CIVIC_2022
export BLOCK=camerad,loggerd,encoderd,micd,logmessaged

# V2V config for the lead role
export HCC_V2V_ENABLED=1
export HCC_V2V_ONLY=0
export HCC_V2V_DEVICE_ID=hcc-lead
export HCC_V2V_RELAY_HOST=127.0.0.1
export HCC_V2V_RELAY_PORT=19090

python -c "
from openpilot.selfdrive.test.helpers import set_params_enabled
set_params_enabled()
from openpilot.common.params import Params
p = Params()
p.put_bool('AlphaLongitudinalEnabled', True)
p.put_bool('EnableHCCC', True)
"

cd system/manager && exec python manager.py
\end{lstlisting}

\subsubsection{Terminal 3 --- Ego openpilot}

\begin{lstlisting}[language=bash]
cd ~/openpilot-hcc
source .venv/bin/activate

export SIMULATION=1 NOBOARD=1 SKIP_FW_QUERY=1
export OPENPILOT_PREFIX=hccego
export FINGERPRINT=HONDA_CIVIC_2022
export BLOCK=camerad,loggerd,encoderd,micd,logmessaged

# V2V config for the ego role
export HCC_V2V_ENABLED=1
export HCC_V2V_ONLY=1
export HCC_V2V_DEVICE_ID=hcc-ego
export HCC_V2V_RELAY_HOST=127.0.0.1
export HCC_V2V_RELAY_PORT=19090

python -c "
from openpilot.selfdrive.test.helpers import set_params_enabled
set_params_enabled()
from openpilot.common.params import Params
p = Params()
p.put_bool('AlphaLongitudinalEnabled', True)
p.put_bool('EnableHCCC', True)
"

cd system/manager && exec python manager.py
\end{lstlisting}

\subsubsection{Terminal 4 --- MetaDrive Bridge}

Wait until both openpilot instances show processes starting, then:

\begin{lstlisting}[language=bash]
cd ~/openpilot-hcc
source .venv/bin/activate
OPENPILOT_PREFIX=hccego python tools/sim/run_bridge.py \
    --mode hc3 --lead_prefix hcclead
\end{lstlisting}

\subsection{Expected Behavior}

\begin{itemize}
  \item MetaDrive shows the ego vehicle following the lead vehicle.
  \item The relay terminal shows periodic traffic (hello + data packets
        in both directions).
  \item The ego's \texttt{controlsState} log will show
        \texttt{hcc\_v2v\_mode\_transition} messages as V2V state
        changes.
  \item With \texttt{HCC\_V2V\_ONLY=1} on the ego, hCCC relies
        exclusively on V2V data (no radar fallback).
\end{itemize}

\subsection{V2V Configuration Reference}

\begin{center}
\begin{tabular}{@{}lll@{}}
\toprule
\textbf{Variable} & \textbf{Ego} & \textbf{Lead} \\
\midrule
\texttt{HCC\_V2V\_ENABLED}    & \texttt{1} & \texttt{1} \\
\texttt{HCC\_V2V\_ONLY}       & \texttt{1} & \texttt{0} \\
\texttt{HCC\_V2V\_DEVICE\_ID} & \texttt{hcc-ego} & \texttt{hcc-lead} \\
\texttt{HCC\_V2V\_RELAY\_HOST}& \texttt{127.0.0.1} & \texttt{127.0.0.1} \\
\texttt{HCC\_V2V\_RELAY\_PORT}& \texttt{19090} & \texttt{19090} \\
\bottomrule
\end{tabular}
\end{center}


% ======================================================================
\section{Mode 3: Hardware-in-the-Loop (HIL)}
\label{sec:mode3}
% ======================================================================

Two Comma~3X devices run the full openpilot stack on-device.  The PC
runs only MetaDrive and ships synthesized camera frames + CAN data to
each device over USB-RNDIS.  V2V communication flows directly between
the devices over WiFi (the PC is not in the V2V path).

\subsection{Hardware Required}

\begin{itemize}
  \item Two Comma~3X devices (one ``ego'', one ``lead'')
  \item Two USB-C cables (PC $\leftrightarrow$ each device)
  \item PC running Linux (recommended) or macOS with HoRNDIS
  \item Optional: Logitech G29/G920 wheel + pedals connected to the PC
\end{itemize}

\subsection{Network Topology}

\begin{center}
\begin{tabular}{@{}lll@{}}
\toprule
\textbf{Link} & \textbf{Transport} & \textbf{Purpose} \\
\midrule
PC $\leftrightarrow$ Ego   & USB-C / RNDIS & Camera, CAN, carControl \\
PC $\leftrightarrow$ Lead  & USB-C / RNDIS & Camera, CAN, carControl \\
Ego $\leftrightarrow$ Lead & WiFi (ego hotspot) & V2V relay only \\
\bottomrule
\end{tabular}
\end{center}

% ── One-time setup ─────────────────────────────────────────────────────
\subsection{One-Time Device Setup}
\label{sec:mode3-onetime}

These steps persist across reboots and only need to be done once per
device.

\subsubsection{Step 1: Flash and Build openpilot on Each Device}

SSH into each device and ensure the \texttt{hcc-ego} branch is checked
out and built:

\begin{lstlisting}[language=bash]
# Repeat for each device
ssh comma@<device-ip>
cd /data/openpilot
git remote add hcc https://github.com/ethanmathias/openpilot-hcc.git
git fetch hcc
git checkout hcc-ego
git lfs pull
git submodule update --init --recursive
scons -j4
\end{lstlisting}

\warn{The \texttt{scons} build on-device takes 15--30 minutes.  A
\texttt{RecursionError} from tinygrad during the build is harmless for
HIL since \texttt{modeld} runs independently and the tinygrad runtime
model is already compiled.}

\subsubsection{Step 2: Configure USB-RNDIS Networking}

Plug both devices into the PC via USB-C.  On the PC, verify the RNDIS
interfaces appear:

\begin{lstlisting}[language=bash]
# Linux
ip link show | grep enx

# macOS (requires HoRNDIS driver)
ifconfig | grep -A2 "en[0-9]"
\end{lstlisting}

Copy and edit the devices configuration:

\begin{lstlisting}[language=bash]
cd ~/openpilot-hcc
cp tools/sim/hil/devices.example.toml tools/sim/hil/devices.toml
\end{lstlisting}

Edit \texttt{tools/sim/hil/devices.toml} with the real interface names
and IPs:

\begin{lstlisting}
[ego]
iface = "enx00e04c360002"    # PC-side interface name
ip = "192.168.32.11"         # Device-side static IP

[lead]
iface = "enx00e04c360001"
ip = "192.168.32.10"
\end{lstlisting}

Set static IPs on each device:

\begin{lstlisting}[language=bash]
# On the EGO device:
ssh comma@<ego-usb-ip>
sudo bash /data/openpilot/tools/sim/hil/scripts/setup_rndis.sh \
    192.168.32.11 192.168.32.1

# On the LEAD device:
ssh comma@<lead-usb-ip>
sudo bash /data/openpilot/tools/sim/hil/scripts/setup_rndis.sh \
    192.168.32.10 192.168.32.1
\end{lstlisting}

On the PC, assign IP addresses to each RNDIS interface:

\begin{lstlisting}[language=bash]
# Linux example:
sudo ip addr add 192.168.32.1/24 dev enx00e04c360002  # ego link
sudo ip addr add 192.168.32.1/24 dev enx00e04c360001  # lead link
\end{lstlisting}

Verify connectivity:

\begin{lstlisting}[language=bash]
ping -c 2 192.168.32.11   # ego device
ping -c 2 192.168.32.10   # lead device
\end{lstlisting}

\subsubsection{Step 3: Configure V2V WiFi Network}

The ego device runs a WiFi hotspot; the lead device joins it.  The V2V
relay server runs on the ego device itself.

\begin{lstlisting}[language=bash]
# On the EGO device:
ssh comma@192.168.32.11
sudo bash /data/openpilot/tools/sim/hil/scripts/setup_v2v_network.sh ego
# Note the hotspot IP printed (typically 10.42.0.1)

# On the LEAD device:
ssh comma@192.168.32.10
sudo bash /data/openpilot/tools/sim/hil/scripts/setup_v2v_network.sh lead
\end{lstlisting}

This script automatically:
\begin{itemize}
  \item Creates and activates the WiFi hotspot (ego) or joins it (lead)
  \item Installs and enables the \texttt{hcc-v2v-relay.service} systemd
        unit on the ego
  \item Sets all V2V params (\texttt{HCCV2VEnabled},
        \texttt{HCCV2VRelayHost}, etc.) on each device
\end{itemize}

Verify from the lead device:

\begin{lstlisting}[language=bash]
ssh comma@192.168.32.10
ping -c 2 10.42.0.1    # should reach ego's hotspot
\end{lstlisting}

Verify the relay is running on the ego:

\begin{lstlisting}[language=bash]
ssh comma@192.168.32.11
systemctl status hcc-v2v-relay
\end{lstlisting}

% ── Running Mode 3 ────────────────────────────────────────────────────
\subsection{Running Mode 3}

You need \textbf{three terminals}.

\subsubsection{Terminal 1 --- Launch the Ego Device}

\begin{lstlisting}[language=bash]
ssh comma@192.168.32.11
bash /data/openpilot/tools/sim/hil/scripts/launch_device.sh \
    ego 192.168.32.1
\end{lstlisting}

\subsubsection{Terminal 2 --- Launch the Lead Device}

\begin{lstlisting}[language=bash]
ssh comma@192.168.32.10
bash /data/openpilot/tools/sim/hil/scripts/launch_device.sh \
    lead 192.168.32.1
\end{lstlisting}

\subsubsection{Terminal 3 --- Launch MetaDrive on the PC}

\begin{lstlisting}[language=bash]
cd ~/openpilot-hcc
source .venv/bin/activate

# With keyboard input:
python tools/sim/run_bridge.py --mode hcc_hil --lead --keyboard

# Or with auto-detected Logitech wheel:
python tools/sim/run_bridge.py --mode hcc_hil --lead
\end{lstlisting}

\subsection{Expected Behavior}

\begin{itemize}
  \item \textbf{PC:} A pygame window appears with three panes---top-down
        view (left), ego camera POV (top-right), lead camera POV
        (bottom-right).  Telemetry overlay shows \texttt{vEgo},
        \texttt{vLead}, \texttt{dRel}.
  \item \textbf{Devices:} The manager starts with
        \texttt{remote\_sensor\_bridge} active.  Native
        \texttt{camerad}, \texttt{pandad}, and \texttt{sensord} are
        blocked.
  \item \textbf{Engagement:} Both devices' \texttt{selfdriveState.active}
        becomes \texttt{True} after the MetaDrive warmup period.
  \item \textbf{Wheel/pedals:} Steering and throttle/brake affect the
        \emph{ego vehicle only}.  The lead vehicle is autonomous.
  \item \textbf{V2V:} Ego's hCCC receives V2V data from the lead via
        the WiFi relay.  The PC is not in the V2V path.
\end{itemize}

\subsection{What \texttt{launch\_device.sh} Does}

For reference, the device launch script sets these environment variables
before starting the openpilot manager:

\begin{center}
\begin{tabular}{@{}ll@{}}
\toprule
\textbf{Variable} & \textbf{Value} \\
\midrule
\texttt{SIMULATION}         & \texttt{1} \\
\texttt{NOBOARD}            & \texttt{1} \\
\texttt{SKIP\_FW\_QUERY}    & \texttt{1} \\
\texttt{HIL\_MODE}          & \texttt{1} \\
\texttt{HIL\_ROLE}          & \texttt{ego} or \texttt{lead} \\
\texttt{HIL\_PC\_IP}        & The PC's RNDIS IP \\
\texttt{BLOCK}              & camerad, pandad, sensord, loggerd, \\
                            & encoderd, micd, logmessaged, ui \\
\bottomrule
\end{tabular}
\end{center}

It also starts two cereal bridge subprocesses per device (msgq$\to$zmq
and zmq$\to$msgq) that relay services between the device and the PC.


% ======================================================================
\section{Troubleshooting}
\label{sec:troubleshooting}
% ======================================================================

\subsection{Common Issues (All Modes)}

\begin{center}
\begin{tabular}{@{}p{5cm}p{9cm}@{}}
\toprule
\textbf{Symptom} & \textbf{Fix} \\
\midrule
\texttt{ImportError: hCCC} &
  The class was renamed to \texttt{HCCC}. Pull latest
  \texttt{hcc-ego} and rebuild. \\
openpilot processes crash on start &
  Ensure \texttt{scons} completed successfully. Check that
  submodules are initialized (\texttt{git submodule update --init}). \\
MetaDrive window is black &
  Check GPU drivers. Try setting
  \texttt{export DISPLAY=:0} on Linux. \\
Scenario replay produces no CSV &
  Verify the Scenarios.csv exists at
  \texttt{tools/sim/lib/Scenarios.csv}. \\
\bottomrule
\end{tabular}
\end{center}

\subsection{Mode 2 Issues}

\begin{center}
\begin{tabular}{@{}p{5cm}p{9cm}@{}}
\toprule
\textbf{Symptom} & \textbf{Fix} \\
\midrule
V2V relay shows no traffic &
  Verify both \texttt{HCC\_V2V\_ENABLED=1} and the relay is running on
  \texttt{127.0.0.1:19090}. \\
Ego ignores V2V, uses radar &
  Set \texttt{HCC\_V2V\_ONLY=1} on the ego instance. \\
Port collision between prefixes &
  Each \texttt{OPENPILOT\_PREFIX} binds its own socket namespace.
  Ensure the two terminals use different prefixes. \\
\bottomrule
\end{tabular}
\end{center}

\subsection{Mode 3 (HIL) Issues}

\begin{center}
\begin{tabular}{@{}p{5cm}p{9cm}@{}}
\toprule
\textbf{Symptom} & \textbf{Fix} \\
\midrule
\texttt{no iframe} spam on device &
  Camera encoder hasn't sent SPS/PPS keyframe. Restart the PC bridge.
  Verify port reachability with \texttt{nc -zv <ip> 19100}. \\
modeld hangs / \texttt{BAD\_FRAME\_ID} &
  Packet loss on USB. Try \texttt{--raw\_yuv} flag on the PC bridge.
  Run \texttt{iperf3} to verify USB bandwidth. \\
\texttt{controlsAllowed=False} &
  \texttt{pandaStates} not arriving on-device. Check cereal bridge logs
  at \texttt{/tmp/hil\_bridge\_*.log}. Verify
  \texttt{OPENPILOT\_PREFIX} matches. \\
Ego ignores wheel input &
  Verify \texttt{user\_torque} is nonzero in the ego's
  \texttt{carState}. The lead should always have zero manual input. \\
V2V not flowing &
  Run \texttt{systemctl status hcc-v2v-relay} on the ego. Try
  \texttt{ping 10.42.0.1} from the lead. Check WiFi hotspot is active
  (\texttt{nmcli device status}). \\
RNDIS interface not appearing &
  Linux: ensure \texttt{cdc\_ether} and \texttt{rndis\_host} kernel
  modules are loaded. macOS: install HoRNDIS. \\
\bottomrule
\end{tabular}
\end{center}


% ======================================================================
\section{Verification Checklist}
\label{sec:verification}
% ======================================================================

Use this checklist to confirm each mode is working correctly before
moving to the next.

\subsection{Mode 1 Checklist}

\begin{enumerate}
  \item[\(\square\)] openpilot starts without import or param errors
  \item[\(\square\)] MetaDrive window opens showing ego behind lead
  \item[\(\square\)] Auto-engagement occurs within 10 seconds
  \item[\(\square\)] Ego speed tracks lead speed (visible in terminal
        status output)
  \item[\(\square\)] \texttt{carControl.actuators.accel} is nonzero when
        engaged
\end{enumerate}

\subsection{Mode 2 Checklist}

\begin{enumerate}
  \item[\(\square\)] All Mode 1 checks pass
  \item[\(\square\)] Relay terminal shows bidirectional packet traffic
  \item[\(\square\)] Ego logs show \texttt{hcc\_v2v\_mode\_transition
        mode=v2v\_active}
  \item[\(\square\)] Ego tracks lead speed even with
        \texttt{HCC\_V2V\_ONLY=1} (radar disabled)
\end{enumerate}

\subsection{Mode 3 Checklist}

\begin{enumerate}
  \item[\(\square\)] \texttt{ping} succeeds from PC to both device RNDIS
        IPs
  \item[\(\square\)] \texttt{ping} succeeds from lead to ego hotspot IP
  \item[\(\square\)] \texttt{systemctl status hcc-v2v-relay} shows
        \texttt{active (running)} on ego
  \item[\(\square\)] Both devices show \texttt{remote\_sensor\_bridge} in
        their process list
  \item[\(\square\)] \texttt{camerad}, \texttt{pandad}, \texttt{sensord}
        are \emph{not} in the process list
  \item[\(\square\)] Pygame window shows three non-black panes
  \item[\(\square\)] Both devices reach
        \texttt{selfdriveState.active = True}
  \item[\(\square\)] Wheel/pedal input deflects the ego only
  \item[\(\square\)] V2V telemetry visible in the ego HUD overlay
\end{enumerate}


% ======================================================================
\section{Quick Reference}
\label{sec:quickref}
% ======================================================================

\subsection{Mode 1 --- Two Commands}

\begin{lstlisting}[language=bash]
# Terminal 1: openpilot
bash tools/sim/launch_openpilot.sh

# Terminal 2: MetaDrive
OPENPILOT_PREFIX=hccego python tools/sim/run_bridge.py --mode hc3
\end{lstlisting}

\subsection{Mode 2 --- Four Commands}

\begin{lstlisting}[language=bash]
# Terminal 1: V2V relay
python tools/hcc_v2v/relay_server.py --host 127.0.0.1 --port 19090

# Terminal 2: lead openpilot
OPENPILOT_PREFIX=hcclead bash tools/sim/launch_openpilot.sh

# Terminal 3: ego openpilot
bash tools/sim/launch_openpilot.sh

# Terminal 4: MetaDrive
OPENPILOT_PREFIX=hccego python tools/sim/run_bridge.py \
    --mode hc3 --lead_prefix hcclead
\end{lstlisting}

\warn{For Mode~2, the lead and ego launch scripts need different V2V
environment variables.  See Section~\ref{sec:mode2} for the full
export list.  The condensed form above works only if
\texttt{launch\_openpilot.sh} is configured to read V2V env vars.}

\subsection{Mode 3 --- Three Commands}

\begin{lstlisting}[language=bash]
# Terminal 1 (SSH to ego device):
bash /data/openpilot/tools/sim/hil/scripts/launch_device.sh \
    ego 192.168.32.1

# Terminal 2 (SSH to lead device):
bash /data/openpilot/tools/sim/hil/scripts/launch_device.sh \
    lead 192.168.32.1

# Terminal 3 (PC):
python tools/sim/run_bridge.py --mode hcc_hil --lead --keyboard
\end{lstlisting}


% ======================================================================
\section{File Reference}
\label{sec:files}
% ======================================================================

\begin{center}
\small
\begin{tabular}{@{}lp{8cm}@{}}
\toprule
\textbf{File} & \textbf{Purpose} \\
\midrule
\multicolumn{2}{@{}l}{\textit{hCCC controller}} \\
\texttt{selfdrive/controls/lib/hccc\_controller.py} & Core HCCC class (feedforward + proportional controller) \\
\texttt{selfdrive/controls/lib/longcontrol.py} & Longitudinal control; owns HCCC instance and V2V subscriber \\
\texttt{selfdrive/controls/controlsd.py} & Top-level controls daemon; wires HCCC into actuator loop \\
\midrule
\multicolumn{2}{@{}l}{\textit{Simulation bridge}} \\
\texttt{tools/sim/run\_bridge.py} & CLI entry point for all bridge modes \\
\texttt{tools/sim/bridge/common.py} & Base \texttt{SimulatorBridge} class and input dispatch \\
\texttt{tools/sim/bridge/metadrive/metadrive\_bridge.py} & MetaDrive-specific bridge (map + scenario config) \\
\texttt{tools/sim/bridge/metadrive/metadrive\_process.py} & MetaDrive subprocess worker (physics, lead vehicle, CSV logging) \\
\texttt{tools/sim/lib/simulated\_car.py} & Simulated Honda Civic CAN message packer \\
\texttt{tools/sim/lib/simulated\_sensors.py} & Simulated IMU, GPS, peripheral, and camera publisher \\
\midrule
\multicolumn{2}{@{}l}{\textit{HIL infrastructure}} \\
\texttt{tools/sim/hil/devices.example.toml} & Template device config (copy to \texttt{devices.toml}) \\
\texttt{tools/sim/hil/device\_config.py} & TOML loader and RNDIS sanity check \\
\texttt{tools/sim/hil/camera\_encoder.py} & H.264 / raw-YUV camera streamer (PC $\to$ device) \\
\texttt{tools/sim/hil/cereal\_bridges.py} & Spawns cereal bridge subprocesses per device \\
\texttt{tools/sim/hil/remote\_sensors.py} & Publishes sensor data to device via ZMQ \\
\texttt{tools/sim/hil/remote\_sensor\_bridge.py} & Device-side ZMQ $\to$ VisionIPC decoder \\
\texttt{tools/sim/hil/window.py} & Pygame composite window (top-down + 2 POVs) \\
\texttt{tools/sim/hil/launch\_pc.py} & PC orchestrator for HIL mode \\
\texttt{tools/sim/hil/scripts/} & Device setup scripts (RNDIS, V2V, launch) \\
\midrule
\multicolumn{2}{@{}l}{\textit{V2V}} \\
\texttt{tools/hcc\_v2v/relay\_server.py} & UDP relay server entry point \\
\texttt{selfdrive/controls/lib/hcc\_v2v.py} & V2V publisher/subscriber and signal types \\
\bottomrule
\end{tabular}
\end{center}

\end{document}
